\section{A uniform bound through the gaps between Fourier levels}
We prove a lower bound for every sample size using a thin middle band
of Fourier modes whose covariance trace is small. Ordinary Lipschitz
contraction controls this band, while Lewis-weight sampling controls
the lower degrees.

\begin{lemma}[A fixed-threshold sampling bound]\label{lem:multiband-process}
Let $r<u<d$, and put $\Delta=\tau_r-\tau_u$. Let
$\mathcal B_{r,u}$ be the span of Walsh characters of degrees $r+1$
through $u$. Suppose
\[
 \eta^2\ge C\frac{D_r\log(2n)}n
\]
for a sufficiently large universal constant $C$. For deterministic
$R,T\ge0$, the class
\[
 \mathcal H=\{-f(p+b):p\in V_r,\ \mathbb E|p|\le R,
                 \ b\in\mathcal B_{r,u},\ \|b\|_K\le\sqrt n\}
\]
satisfies
\begin{equation}\label{eq:multiband-process}
 \mathbb E_S\sup_{h\in\mathcal H}
   |\mathbb P_n(h-T)_+-\mathbb E_x(h-T)_+|
 \le4(R\eta+\sqrt\Delta).
\end{equation}
Here $\mathbb P_n g=n^{-1}\sum_i g(X_i)$.
\end{lemma}
\begin{proof}
For the full-cube matrix with rows $f(x)(\chi_U(x))_{|U|\le r}$,
the Lewis weights are $D_r/M$. The linear-process moment bound of
Mai, Musco, and Rao \cite[Theorem 3]{mai2021coresets}, with its failure
parameter set to $1/2$, therefore implies
\[
 \mathbb E_{S,\epsilon}\sup_{p\in V_r,\,\mathbb E|p|\le R}
 \left|\frac1n\sum_i\epsilon_i f(X_i)p(X_i)\right|\le R\eta.
\]
The normalization follows directly: their sampled matrix rescales each
sampled row by $M/n$, whereas $\|p\|_1=M^{-1}\|A\beta\|_1$.
Their moment bound implies the displayed expectation bound by Jensen.

The covariance-weighted feature vector of the band has squared norm
$\Delta$ at every input. Consequently
\[
 \mathbb E_{S,\epsilon}\sup_{\|b\|_K\le\sqrt n}
 \left|\frac1n\sum_i\epsilon_i f(X_i)b(X_i)\right|
 \le\frac{\sqrt n}{n}\sqrt{n\Delta}=\sqrt\Delta.
\]
The map $z\mapsto(z-T)_+$ is $1$-Lipschitz and vanishes at zero.
Symmetrization, followed by the scalar contraction principle, bounds
the left-hand side of \eqref{eq:multiband-process} by four times the
linear Rademacher width of $\mathcal H$. The triangle inequality and
the preceding two estimates prove the claim. Both parameter balls and
the threshold are fixed before sampling; no positive homogeneity in
$T$ or relative bound on an unbounded parameter class is asserted.
\end{proof}

\begin{theorem}[Every sample size up to an exponential threshold]
\label{thm:all-n-gaussian}
Fix $L\ge2$. There exist $c_L'>0$, $\varepsilon>0$, and $d_0$ such
that for every $d\ge d_0$ and every integer $0\le n\le e^{c_L'd}$,
independent uniform parity training inputs satisfy
\[
 \mathbb P_S\{F_{13/20}(S)\ge\varepsilon\}\ge\frac34.
\]
The constants may be chosen with
$\varepsilon=c_\theta^2/256$, where
$c_\theta=\sqrt{2\pi}(1/e-7/20)$. The same conclusion holds for
$\liminf_{N\to\infty}F_N(S)$ at every fixed $\gamma<0$ in the
network convention of Proposition~\ref{prop:paper-width}.
\end{theorem}

We first prove the range $e^{(\log(d+1))^3}\le n\le e^{c_L'd}$.
A short argument for smaller samples follows the proof.

\begin{proof}[Proof in the larger-sample range]
Fix a small sampling-failure probability $\delta=1/20$. All constants
below may depend on $L$ and $\delta$. Choose a sufficiently large
constant $A$, to be fixed below, and let $r$ be the largest odd integer
such that
\begin{equation}\label{eq:coarse-cutoff}
 A D_r\log(D_r+1)\le n.
\end{equation}
By choosing $c_L'$ sufficiently small, for all large $d$ this cutoff
satisfies
\[
 \tfrac14(\log d)^2\le r\le\alpha d,
 \qquad n<A D_{r+2}\log(D_{r+2}+1),
\]
where $\alpha>0$ is a small constant chosen after $A$. Indeed
$D_r\le(ed/r)^r$ gives the lower bound, while binomial entropy
estimates give the upper bound. Uniformly in this range,
$\log(2n)\le C\log(D_r+1)\le C r\log(eD/r)$.

Put $\tau=\tau_r$, $\ell=\lambda_{r+1}$, $s_2=s_{2,r}$, and
\[
 B=C_\delta\left(1+\frac{\sqrt{nd s_2}}\tau+
                            \frac{n\ell}\tau\right),\qquad
 \eta^2=C_0\frac{D_r\log(2n)}n.
\]
The constant $C_0$ is large enough for
Lemma~\ref{lem:multiband-process}. The random Gram estimate implies
$\|Q_{r,SS}\|\le B\tau$ with probability at least $1-\delta$ for
all large $d$. Its omitted far-pair terms tend to zero uniformly here:
$r\ge(\log d)^2/4$, and $c_L'$ is small enough that
$2c_L'<\log2-h(1/4)$.

The crucial estimate is
\begin{equation}\label{eq:variable-accuracy}
 \eta^2B\le C_\delta\left[
 \frac1A+\sqrt{\frac{\rho\log(e/\rho)}A}
                +\rho^2\log(e/\rho)\right],\qquad \rho=r/D.
\end{equation}
To verify it, first $\eta^2\le C/A$ by \eqref{eq:coarse-cutoff}.
Next, the sharp spectral estimates give
\[
 \eta^2\frac{n\ell}\tau
 \le\frac{C D_r\log(2n)}{r\binom D{r+2}}
 \le C\rho^2\log(e/\rho).
\]
Finally, the square of the remaining term is at most
\[
 \left(\eta^2\frac{\sqrt{nd s_2}}\tau\right)^2
 \le\frac CA\frac{C D_r\log(2n)\,d}{r^2\binom D{r+2}}
 \le\frac CA\rho\log(e/\rho).
\]
This proves \eqref{eq:variable-accuracy}. Its right-hand side can be
made arbitrarily small, uniformly over all these sample sizes, by
choosing $A$ large and then $\alpha$ small.

Let $m$ be the actual sign-posterior mean. Its degree-$r$ truncation
$p$ lies, on the Gram and ordinary ReLU sampling events, in the
deterministic $L^1$ ball
\begin{equation}\label{eq:head-radius}
 \mathbb E|p|\le4\sqrt{B\tau}=:R.
\end{equation}
Indeed $\|m-p\|_K^2\le n$ and the Gram event bound its empirical
squared norm by $B\tau$. Training consistency and the ordinary ReLU
estimate at accuracy $1/4$ then give \eqref{eq:head-radius}, without
a VC or squared-loss estimate. For sufficiently large $A$, that ReLU
event has probability at least $1-\delta$.

Now choose the odd cutoff
\[
 u=r+2\lceil10\log(d+1)\rceil,
\]
and split $m=p+b+t$ into degrees at most $r$, degrees $r+1$ through
$u$, and degrees greater than $u$. The sharp level masses give
\begin{equation}\label{eq:thin-band}
 \Delta:=\tau_r-\tau_u\le C\tau_r\frac{\log d}{r}=o(\tau_r)
\end{equation}
uniformly in the sample range. The upper shell threshold bounds
$n\ell/\tau\le CA\log(eD/r)$ and
$nd s_2/\tau^2\le CA(d/r)\log(eD/r)$.
Below degree $D/8$, each successive odd binomial coefficient grows by
at least a factor four. The additional $\lceil10\log(d+1)\rceil$
steps therefore ensure
\[
 n\lambda_{u+1}/\tau=o(1),\qquad
 nd s_{2,u}/\tau^2=o(1),\qquad \tau_u/\tau=1-o(1).
\]
Apply the random Gram estimate at $u$. With probability $1-o(1)$,
\begin{equation}\label{eq:final-tail}
 \mathbb P_n t^2\le2\tau,\qquad
 \|t\|_2^2\le n\lambda_{u+1}=o(\tau).
\end{equation}
The normalized trace of the tail covariance captured by training is
at most $n\lambda_{u+1}$. The residual-variance argument thus also
gives $\sigma_x^2\ge\tau/2$ outside an $o(1)$ fraction of inputs.
All these error terms are uniform in $n$.

Set $a_0=c_\theta/(2\sqrt2)$ and choose the deterministic threshold
$T=4\sqrt\tau/a_0$. Lemma~\ref{lem:multiband-process}, Markov's
inequality, and \eqref{eq:head-radius} imply, with probability at
least $1-\delta$, a uniform error at most
\[
 \frac4\delta(R\eta+\sqrt\Delta)
\]
for the empirical stoploss of every member of the indicated sum class.
By \eqref{eq:variable-accuracy} and \eqref{eq:thin-band}, choose $A$
large, $\alpha$ small, and then $d$ large so that this is at most
$a_0\sqrt\tau/8$. The actual $p+b$ belongs to this class on the
events already specified, since $\|b\|_K\le\sqrt n$.

Write $h=p+b$ and $F=F_{13/20}(S)$. The confidence margin and
\eqref{eq:final-tail} imply
$fh>a_0\sqrt\tau$ on a set of measure at least $1-F-o(1)$.
Its negative set consequently has measure at most $F+o(1)$.
Moreover $\mathbb E fh=0$ exactly, because $u<d$.
Training consistency gives $(-f(X_i)h(X_i))_+\le|t(X_i)|$, and
$(z-T)_+\le z^2/(4T)$ for $z\ge0$. Hence
\[
 \mathbb P_n(-fh-T)_+\le\frac{\mathbb P_nt^2}{4T}
             \le\frac{a_0\sqrt\tau}{8}.
\]
The uniform stoploss estimate raises this bound to at most
$a_0\sqrt\tau/4$ in the population. Equality of the positive and
negative integrals of $fh$ now yields
\[
 a_0\sqrt\tau(1-F-o(1))
 \le\mathbb E(-fh)_+
 \le T(F+o(1))+\frac{a_0\sqrt\tau}{4}.
\]
After division by $\sqrt\tau$, this implies
$F\ge a_0^2/32=c_\theta^2/256$ for all large $d$.
The three events with failure probability $\delta$, together with
the event of failure probability $o(1)$, have joint probability
at least $1-3\delta-o(1)>3/4$. This completes the larger-sample proof.
\end{proof}

\begin{lemma}[Parity obstruction for polynomial thresholds]
\label{lem:small-ptf}
Every polynomial threshold of degree at most $s<d$ has full-parity
error at least
\[
 B(d,s)=2^{-d}\sum_{j=0}^{\lfloor(d-s-1)/2\rfloor}\binom dj.
\]
\end{lemma}
\begin{proof}
A sufficiently small constant perturbation removes zero ties without
changing the threshold's assigned signs. Let $P$ be the resulting
polynomial and $t=\lfloor(d-s-1)/2\rfloor$. If the error set has fewer
than $\dim V_t$ points, linear algebra supplies a nonzero $Q\in V_t$
vanishing at every error. Then $fPQ^2$ is everywhere nonnegative and
strictly positive somewhere. But $\deg(PQ^2)\le s+2t<d$, so its
correlation with full parity is zero, a contradiction.
\end{proof}

\begin{proof}[Completion for smaller samples]
For $n\le e^{(\log(d+1))^3}$ choose an odd $s\sim d^{1/3}$. The
sharp spectral bound gives $n\lambda_{s+1}/\tau_s=o(1)$, uniformly
in $n$. The posterior mean norm bound, residual-variance lemma, and
confidence margin imply that its degree-$s$ truncation agrees with
parity on every confident input except an $o(1)$ fraction of inputs.
Thus Lemma~\ref{lem:small-ptf} gives
\[
 F_{13/20}(S)\ge B(d,s)-o(1)\ge\frac13
\]
for all large $d$, for every training sample. Here $B(d,s)\to1/2$:
by symmetry the removed central binomial window has $O(s)$ atoms,
each of mass $O(d^{-1/2})$, and $s=o(\sqrt d)$.
This proves the remaining range. Proposition~\ref{prop:paper-width}
transfers the confidence conclusion to every fixed negative exponent.
\end{proof}
